% ===========================================================================
\section{Benchmarks, Including Two Disagreements Worth Defending}
\label{sec:benchmarks}
% ===========================================================================

Every sector is compared against whichever external estimates exist for it: EDGAR
\citep{edgar2025booklet} for all nine non-agricultural sectors, FAOSTAT \citep{faostat} and Climate
TRACE \citep{climatetrace} for agriculture, and Pakistan's own BUR-1 \citep{pakbur1} where its
categories are separable. The comparison is generated on every run rather than compiled by
hand, and \num{26} of its \num{27} rows resolve against an engine category; the one that does
not is the red-flagged CRT ammonia comparator, which reports native \COtwo{} and is never
divided into a \COe{} value. Each row's register entry declares its GWP basis, and the engine
side is resolved in that basis: Climate TRACE publishes under AR6, so those seven rows compare
AR6 against AR6, where an earlier edition compared them against AR5 values and manufactured a
few percent of disagreement in the comparison itself. Resolution does not
imply completeness: the industrial row measures the metered gas component against EDGAR's
all-fuel category, so its low ratio remains visible.

\begin{table*}[htbp]
\centering
\footnotesize
\setlength{\tabcolsep}{5pt}
\begin{tabular}{@{}p{0.26\linewidth}llrrr@{}}
\toprule
\textbf{Comparison} & \textbf{Source} & \textbf{Period} & \textbf{Bench.} & \textbf{Engine} &
\textbf{Ratio} \\
\midrule
\multicolumn{6}{@{}l}{\emph{Comparisons on matched boundaries}} \\
Electricity (1A1a) & EDGAR & CY2024 & 56.80 & 40.04 & 0.70 \\
Residential (households) & EDGAR & CY2024 & 31.42 & 24.22 & 0.77 \\
Waste (year-matched) & EDGAR & CY2023 & 26.57 & 37.45 & 1.41 \\
Agriculture total & FAOSTAT & CY2023 & 233.30 & 223.45 & 0.96 \\
Agriculture total & EDGAR & CY2024 & 261.20 & 223.45 & 0.86 \\
Agriculture total & Climate TRACE & CY2024 & 261.50 & 218.07 & 0.83 \\
Managed soils & FAOSTAT & CY2023 & 58.50 & 28.53 & 0.49 \\
Rice cultivation & Climate TRACE & CY2024 & 7.80 & 15.96 & 2.05 \\
[3pt]
\multicolumn{6}{@{}l}{\emph{Partial-scope diagnostics --- ratio withheld or scope-noted}} \\
Fugitives (1B2 vs all 1B)$^{e}$ & EDGAR & CY2024 & 17.54 & 23.50 & 1.34 \\
Transport module vs EDGAR modes & EDGAR & CY2024 & 48.76 & 41.93 & --- \\
[3pt]
\multicolumn{6}{@{}l}{\emph{Coverage comparisons --- not agreement}} \\
All sectors & EDGAR & CY2024 & 514.33 & 444.96 & 0.87 \\
Industrial combustion (1A2) & EDGAR & CY2024 & 43.98 & 28.13 & 0.64 \\
IPPU: cement term vs all processes & EDGAR & CY2024 & 27.38 & 21.82 & --- \\
\bottomrule
\end{tabular}
\caption{The 1.A.1.a grid figure is the \texttt{power\_grid\_control} family reference; the family's other members read 38.45~Mt (fuel property), 44.05~Mt (coverage) and 57.82~Mt (legacy account), and none is a sensitivity.
$^{e}$~The 1.B.2 reference sits one membership step below the top of its own
9.601--24.694~Mt structural envelope, fertiliser feedstock being outside the post-meter
leakage base; the priced envelope members, including the 17.68~Mt geometric-mean
central reading, sit beside it in the configurations file.
External comparison, \si{\mega\tonne} \COe{}, \textbf{grouped by what the comparison
can support}. A previous edition printed all three groups as one undifferentiated ratio
column, which invited a scope difference to be read as an error---and one row was worse than
that: it was labelled ``Road transport'' while carrying the transport \emph{module}, road plus
mixed off-road, against an EDGAR row for a different set of modes in a different year. Every
row now carries its period, and the register behind it
(\texttt{data/benchmarks.csv}) carries both boundaries and a \texttt{ratio\_permitted} flag; a
dash means the boundaries do not support a ratio, not that one could not be computed. The
inventory period is FY2023--24 throughout, so a row whose benchmark period is a calendar year
is comparing across a partial-year offset as well.}
\label{tab:benchmarks}
\end{table*}

\subsection{The aggregate reading}

At the whole-inventory level this work is \SI{87}{\percent} of EDGAR's Pakistani sectors
summed. That number is a \emph{magnitude ratio} and is reported as one, because on its own it
conflates two questions that different work answers. Source domains this inventory does not
estimate at all pull it down; domains it does estimate and where the two disagree pull it in
either direction---waste sits at \SI{141}{\percent} of EDGAR and agriculture at
\SI{86}{\percent}---and netting those against each other produces a single figure that reads
as completeness and is not one. Three numbers are therefore reported separately
(\texttt{outputs/national\_benchmark\_scope\_coverage.csv}):

\begin{enumerate}[leftmargin=*,itemsep=2pt,topsep=3pt]
\item \textbf{Magnitude ratio, \SI{87}{\percent}}---\SI{444.96}{} against
  \SI{514.33}{\mega\tonne}. How big, not how complete.
\item \textbf{Benchmark-scope representation, \SI{11}{\percent} fully and
  \SI{89}{\percent} partially}---the share of EDGAR's source domain, weighted by EDGAR's own
  magnitudes, in which \emph{every} IPCC category mapped to the domain is estimated, included
  elsewhere or does not occur. Only public electricity meets that. Nothing is wholly absent:
  every domain carries at least one estimated category, which the magnitude ratio cannot show.
  The test is deliberately strict---one unestimated category makes a domain partial, however
  small it is---so the unestimated categories are named per domain rather than absorbed:
  liming and other carbon-containing fertilisers in agriculture, domestic navigation in
  transport, industrial solid fuels in 1.A.2, commercial buildings, four IPPU categories, and
  industrial wastewater, incineration and biological treatment in waste.
  \emph{An earlier edition of this paper claimed five domains estimated ``whole'', which the
  completeness matrix in the same repository contradicted.} The classification is derived from
  the notation keys now, through \texttt{data/edgar\_domain\_map.csv}, rather than asserted.
\item \textbf{Inventory completeness}---the IPCC notation-key matrix over all \num{138}
  category--gas combinations in Section~\ref{sec:completeness}, which is the only one of the
  three that is independent of a benchmark's own scope.
\end{enumerate}

\textbf{The bridge from this inventory to EDGAR's total has three terms, and mixing them is
the error to avoid.} Categories absent here, valued on EDGAR's own basis, are a
\emph{coverage} gap: industrial combustion (\SI{15.85}{\mega\tonne} of its
\SI{43.98}{\mega\tonne}, after the brick and cement coal activation) and F-gases
(\SI{11.54}{\mega\tonne}) are the large ones.
Differences where both systems estimate the same category are a \emph{method or activity}
disagreement, and they run in both directions: agriculture is \SI{38}{\mega\tonne} below
EDGAR and waste is above it. A previous edition listed the agriculture difference beside the
two absences as though it were a third gap, which it is not---nothing is missing there, the
two estimates disagree. Adding only the coverage term at EDGAR's own values would put this
inventory slightly \emph{above} EDGAR overall, which is the
right sanity check on the estimated categories: they are not systematically low,
the inventory remains incomplete in specific places, and the completeness table
(Section~\ref{sec:completeness}) names exactly where.

The resolved sector comparisons fall broadly within a factor of two, which for
independent Tier~1 estimates built from different activity data is close agreement. The
flagged component rows are examined rather than adjusted toward their benchmarks.

\subsection{Railways against Pakistan's own submission: a factor of ten}

BUR-1 reports railways at \SI{4.06}{\mega\tonne} \COe{} for FY2017--18, and the comparison
this paper makes is \emph{year-matched}: run on the operator's metered fuel for that same
fiscal year, by the identical method, this inventory reconstructs \SI{0.4663}{\mega\tonne}.
BUR-1 is \num{8.7} times it. Quoting the current year against a 2017--18 filing would fold six
years of activity change into a ratio that is about method, which is why the matched year
leads; both years are set out together so neither can be read as the other.

\begin{table}[htbp]
\centering
\footnotesize
\begin{tabular}{@{}lrr@{}}
\toprule
\textbf{Quantity} & \textbf{FY2017--18} & \textbf{FY2023--24} \\
\midrule
Reported diesel, thousand litres & \num{156396} & \num{132588} \\
Estimated diesel mass, \si{\kilo\tonne} & \num{132} & \num{112} \\
This inventory, \si{\mega\tonne} \COe{} & \num{0.466} & \num{0.394} \\
BUR-1, \si{\mega\tonne} \COe{} & \num{4.06} & --- \\
Submission over this inventory & \num{8.7} & --- \\
\bottomrule
\end{tabular}
\caption{The railway comparison, matched. \emph{A previous edition mixed three years in one
paragraph}: it quoted the matched-year estimate beside FY2017's \num{144631} thousand litres,
which is a different year again, and beside the current year's mass. The matched-year estimate
is built on the \num{156396} thousand litres the operator reports for FY2017--18. The
difference between the two years is fuel, not method: FY2017--18 diesel is \SI{18.0}{\percent}
above FY2023--24, not a factor of ten in either direction.}
\label{tab:railmatched}
\end{table}

The disagreement resolves on arithmetic rather than on judgement. Producing
\SI{4.06}{\mega\tonne} of \COtwo{} from diesel requires roughly \SI{1280}{\kilo\tonne} of
fuel. Pakistan Railways---a single state operator that meters and publishes its own fuel
purchases---reports buying \SI{112}{\kilo\tonne} in FY2023--24 and \SI{132}{\kilo\tonne} in
FY2017--18, so BUR-1's implied fuel is roughly ten times either. The traffic comparison is a
\emph{physical-scale diagnostic} rather than a matched-year test, and is labelled as one
because it pairs the current year's train-kilometres with the older submission: the metered
purchase works out at \SI{4.0}{\litre} per train-kilometre over the operator's own
\num{33.1} million train-km, which is where a diesel locomotive belongs, while BUR-1's implied
fuel would need about \SI{46}{\litre} per train-kilometre, which no locomotive does. The
conclusion does not rest on it---the fuel-mass route above is matched on both years.

Three independent quantities---the operator's fuel account, the published traffic statistics
and the carbon content of diesel---all support the lower figure, and both the fuel bound and
the per-train-kilometre intensity are pinned by tests so the claim is rechecked on every run.
We record the disagreement openly, on the principle that a national inventory which silently
diverges from its country's own submission is worse than one that diverges and says so.

Three further facts from the same register
(\texttt{railways\_official\_comparison\_2024.csv}) place that disagreement, and none of them
flatters this inventory by construction.

\emph{It is not confined to the superseded narrative filing.} Pakistan's first Common
Reporting Tables, the current submission and the one this paper prefers wherever a cell
exists, book \SI{56609}{\tera\joule} under 1.A.3.c for calendar 2022. Against this
inventory's value for the two fiscal years that calendar year straddles, the filing is
\num{11.6} times it---worse than BUR-1's \num{8.7}, not better.

\emph{It runs the other way against Pakistan's oldest filing.} The Initial National
Communication implies \SI{0.5575}{\mega\tonne} for FY1993--94, which is
\SI{11.9}{\percent} \emph{below} this inventory's independent reconstruction of the same
period from the operator's own fuel history. That figure is derived from a stated share
rather than read from a printed row, and the register labels it a magnitude check rather than
a point match; it is reported here because it is the one official comparison this inventory
does not win, and because of what it implies about timing. Between the only two filings that
report railways separately the official figure rose by a factor of \num{7.3} while this
inventory's reconstruction \emph{fell} by \SI{26}{\percent}. The order-of-magnitude
disagreement is therefore not a standing feature of Pakistan's reporting: it entered the
official record somewhere between 2003 and 2022. It cannot be dated more closely than that,
because the two filings in between report only the 1.A.3 aggregate.

\emph{And it exceeds anything the railway has ever emitted.} The operator's fuel history,
reconstructed back to 1950, peaks at \SI{2.3049}{\mega\tonne} \COe{} per year in
1965--70, when Pakistan Railways burned \num{48720} tonnes of coal and \num{542078} tonnes
of furnace oil annually at the end of the steam era. BUR-1's \SI{4.06}{\mega\tonne} is
\num{1.8} times that peak. A diesel-only railway a quarter the size of the steam network
cannot emit twice what the steam network did, and no assumption about factors or boundaries
recovers the difference---only the activity can, which is where the filing and the operator's
metered purchases part company.

The railway fuel account also supplies this inventory's only genuinely \emph{independent}
source-boundary check, and it is worth separating from the closures of
Section~\ref{sec:closures} because it is the opposite kind of evidence. A closure is an
identity: one quantity computed once and consumed twice, which cannot disagree because it is
the same number. This is two unrelated reporting chains describing the same physical fuel.
Pakistan Railways meters and publishes its own purchases in its Year Book; the Oil Companies
Advisory Council independently reports a Railways customer-class row in its national
petroleum statistics. Neither is derived from the other. Year-matched for FY2021--22 they
differ by \SI{3.80}{\percent}---\SI{112037}{\tonne} against \SI{107932}{\tonne}---which is
about as close as two independent national fuel accounts get.

That check exists because a closure stopped being necessary. The transport sector once
subtracted the railway fuel account from a single national diesel supply figure; since it was
re-anchored on the petroleum statistics' own customer classes, Railways sits outside the Road
row by construction and there is nothing left to subtract. The quantity that used to enforce
a boundary now tests one instead.

\subsection{Waste against EDGAR: an unresolved methane correction factor}

The waste sector is \SI{141}{\percent} of EDGAR's year-matched Pakistani waste total. This row has changed
sides during the project: before domestic wastewater was built it read \SI{47}{\percent}, and
the gap was a measure of what 4D would add. With 4D in, the sector overshoots.

The likeliest single cause is the methane correction factor for sewers. IPCC gives
\num{0.5} for stagnant sewers and \num{0.1} for wastewater discharged to a flowing river. At
\num{0.5}, domestic wastewater is \SI{24.46}{\mega\tonne} and the waste sector overshoots
EDGAR. At \num{0.1} it is \SI{15.39}{\mega\tonne} and the sector lands close.

The evidence points both ways, and genuinely so. The physical description of Pakistani urban
drains---open, slow-moving, heavily organically loaded---supports \num{0.5}. The only
independent comparator supports \num{0.1}. EDGAR is itself a model and does not settle the
question.

We publish both readings rather than resolving the tension by fiat, and report the
\num{0.5} case as the headline because it is the value the physical description supports.
This is the same treatment given to the correlated-uncertainty case in
Section~\ref{sec:uncertainty}: where an assumption moves a result by a factor approaching
two and the evidence does not choose, the defensible output is both numbers and a statement
of what would settle it. Here, what would settle it is a dissolved-methane measurement
campaign on urban drains in two or three Pakistani cities.

\subsection{The other flagged rows}

\textbf{Fugitives at \num{1.34} of EDGAR} is a matched-year but only partly matched-scope
comparison. The 2019 segment account includes production, gathering, processing,
transmission, distribution, refining and post-meter sources; EDGAR labels its aggregate
``fuel exploitation''. The difference lies inside the engine's declared structural and
parametric ranges, but no adjustment is made toward EDGAR because both products depend on
modelled factors rather than Pakistani component measurements.

\textbf{Managed soils at \num{0.49} of FAOSTAT} and \textbf{rice at \num{2.05} of Climate
TRACE} are agriculture-sector questions treated in the companion manuscript
\citep{ahmad2026ag}. Both survive scrutiny, and the managed-soils one has been measured
rather than asserted since this release.
\texttt{outputs/pakemis\_faostat\_nitrogen\_flows.csv} compares the two compilations at the
nitrogen flows rather than at the total, and the flows agree. FAOSTAT's four managed-soils
items map one-to-one onto this inventory's four pools, so none of the difference is scope.
Synthetic fertiliser nitrogen agrees to \SI{0.9}{\percent}---\num{3775000} against
\num{3808775}~t~N---and the six livestock species both compilations carry agree to
\SI{0.9}{\percent} in total, with the residual running the wrong way for the headline: this
inventory counts \num{1.36} times FAOSTAT's sheep and \num{1.30} times its camels, so more
excreted nitrogen, not less.

The factor of two is therefore in the emission factors and in the partition of nitrogen, not
in the nitrogen. Two-fifths of it is the factor vintage: FAOSTAT applies the 2006
Guidelines' aggregate EF1 of \num{0.010} where this inventory applies the 2019 Refinement's
\num{0.005} for dry climates, and taking FAOSTAT's vintage as a named scenario moves managed
soils from \SI{28.53}{\mega\tonne} to \SI{40.43}{\mega\tonne}, closing
\SI{11.90}{\mega\tonne} of a \SI{29.97}{\mega\tonne} gap. The remaining
\SI{18.07}{\mega\tonne} is neither the nitrogen nor that vintage, and it is not diffuse
either. FAOSTAT's four items, acquired separately for this release, put two-thirds of the
whole gap in one pool: pasture deposition is \SI{7.71}{\mega\tonne} here against
\SI{27.29}{\mega\tonne}, a difference of \SI{19.58}{\mega\tonne} and \SI{65.3}{\percent}
of the total. Synthetic fertiliser contributes \SI{8.77}{\mega\tonne}
(\SI{29.3}{\percent}) and crop residues \SI{2.70}{\mega\tonne}, and the mechanism in all
three is the one already named: the 2019 Refinement disaggregated these direct factors by
climate and Pakistan sits at the dry end. Beside EF1's \num{0.005} against \num{0.010}, the
pasture factor for cattle is \num{0.002} here against the 2006 default's \num{0.02}---a
tenfold cut, applied to the \SI{40.9}{\percent} of this inventory's
\SI{7.7834}{\mega\tonne} of excreted nitrogen that its ledger routes to grazing.

Taken together the two direct factors account for almost the whole difference. The pasture
factor had no scenario beside it until this release, which was an asymmetry rather than a
judgement: the same climate choice on EF1 has carried three alternatives for releases---the
wet value, the 2006 aggregate and an empirical adjudicator---while the larger lever carried
none. It has one now, and applying both 2006 direct factors together moves managed soils to
\SI{54.72}{\mega\tonne}, closing \SI{26.19}{\mega\tonne} or \SI{87}{\percent} of the gap.
The two scenarios are disjoint by construction: EF1 scales synthetic fertiliser, organic
manure and crop residues, the pasture factor scales the two deposition rows, and no subitem
belongs to both. What remains is \SI{3.78}{\mega\tonne}, in the indirect factors---EF4 and
EF5 changed between the vintages too, and these scenarios do not touch them---and in the
nitrogen partition.

One pool runs the other way and the factor vintage cannot explain it. Manure applied to soils
is \num{1.16} times FAOSTAT's despite the lower factor, so this inventory routes more
excreted nitrogen to field application than FAOSTAT does. That is a disagreement about where
the nitrogen goes rather than about what happens to it, and it is the only one of the four.

Whether the dry classification is right for grazing land is a separate question from whether
it is right for cropland, and the evidence behind it is weaker. The dry EF1 is recorded as
\emph{best-supported for irrigated arid cropland}, derived largely from irrigated dry
fields---which is what Pakistan's fertilised area is. Pasture, range and paddock land is
neither irrigated nor cropland, so that basis does not transfer to it, and no comparable
Pakistani evidence stands behind the pasture choice. The factor stays at the Refinement's dry
value, because that is the Refinement's own classification for an arid country, and the
alternative is now priced beside it rather than left implicit. Climate TRACE's rice estimate, the other agriculture question
here, uses a cultivated-area basis that diverges from the national crop statistics.

\textbf{IPPU at \num{0.93} of EDGAR} is nearer like-for-like than the label it used to carry,
and what remains unmatched is now named rather than assumed. EDGAR's ``Processes'' row covers
cement, lime, soda ash, ammonia, and iron and steel. Of those five this inventory estimates
three---cement at \SI{21.82}{\mega\tonne}, ammonia at \SI{3.24}{\mega\tonne}, and iron and steel
at \SI{0.39}{\mega\tonne}---and adds lubricant use, \SI{0.08}{\mega\tonne}, which EDGAR's row
does not name. Lime is NE for want of activity data; soda ash is NA rather than NE, its
Solvay-process carbon allocated to Energy under IPCC Volume~3 \S3.8.3.1
(Section~\ref{sec:completeness}). The \SI{1.84}{\mega\tonne} by which EDGAR's row exceeds this
module is therefore of the order of one absent category, not of four, and the comparison is
recorded as a coverage reading rather than as agreement.

Cement alone is \num{0.80} of the whole-processes row, the right order for a country whose IPPU
is dominated by cement. The cement term's own configuration range brackets matched-period
clinker trade and the clinker-fraction choice: the Tier~1 diagnostic at the 0.75 fraction gives
\SI{17.39}{\mega\tonne}, while the 0.95-fraction trade-corrected configuration gives
\SI{21.50}{\mega\tonne}; the reference itself is the producer census's observed clinker,
\SI{21.82}{\mega\tonne}, above both. The earlier 48.011-Mt cement control was resolved as FY2021--22,
not an alternative FY2023--24 series.

\subsection{The national submission, year-matched: the gap is completeness, not estimation}
\label{sec:btr}

The comparators above all require an approximation. EDGAR and Climate TRACE publish calendar
years against this inventory's fiscal one; Climate TRACE publishes under AR6 against an AR5
account; FAOSTAT does both. One comparator requires neither, and it is the most authoritative
one available: Pakistan's First Biennial Transparency Report \citep{pakbtr2024} publishes a
full IPCC-category inventory on AR5 global warming potentials, and it states its own period
convention---``the fiscal year 1993--94 (referred to as 1994)'', ``Pakistan's 2020--21
inventory (referred to as 2021)''. Its inventory year is a Pakistani fiscal year. It therefore
faces FY2020--21 directly, with no blend and no metric conversion, and this release carries
that year.

Four of this project's own ingests had read that document and three of them recorded it as
calendar 2021; only the aviation ingest read it as fiscal. The correction is why this
comparison is possible at all, and it is recorded in the register rather than in a footnote.

A join on the IPCC code would be wrong, and silently so. The submission books rice cultivation
at 3.D and managed soils at 3.C; this inventory books rice at 3.C and soils at 3.D. The codes
are \emph{swapped} between the two accounts, so a code join compares rice against soil
nitrous oxide and returns two plausible ratios. Every row of
\texttt{outputs/pakemis\_btr\_comparison.csv} is therefore an explicit crosswalk entry with a
declared tier, and a ratio is emitted only where the boundaries support one.

Excluding the submission's land-use block---\SI{34.82}{\mega\tonne}, outside this inventory by
declared scope---the two national totals are \SI{483.45}{} and \SI{486.65}{\mega\tonne}, a
ratio of \num{0.993}. That agreement is not the result worth reporting, because it averages
category ratios spanning more than an order of magnitude: \num{6.91} on oil and gas fugitives,
\num{0.44} on managed soils, \num{1.19} on enteric fermentation, \num{0.43} on lubricants.
What the headline does support is a decomposition the single ratio hides. Restricted to the
categories both accounts actually estimate, the totals are \SI{482.49}{} and
\SI{480.61}{\mega\tonne}---a ratio of \num{1.004}, this inventory above the submission.
The \SI{3.20}{\mega\tonne} national difference is the net of two movements in opposite
directions. \SI{6.04}{\mega\tonne} is seven categories this inventory declares NE and the
submission books---solid-fuel fugitives, lime, glass, other carbonate uses, other
industrial processes, waste incineration and industrial wastewater---against
\SI{0.382}{\mega\tonne} this inventory estimates where the submission books nothing or
zero, which is iron and steel: a net \SI{5.658}{\mega\tonne} of pure completeness
difference in the submission's favour. On the categories both accounts estimate, this
inventory is \SI{1.88}{\mega\tonne} above the submission, which offsets part of it.
Those seven rows are the same rows that sit in
\texttt{outputs/pakemis\_completeness\_mass.csv} with no bound, and they now have an
official magnitude.

\textbf{The sign of that national difference reversed with the coal re-base and has
reversed back in this release, and earlier
editions of this paragraph read it each way.} The first described the gap as this inventory's
\emph{shortfall} and concluded that the difference against Pakistan's own account was a
completeness difference rather than an estimation one. The matched comparison year is
inventory year 2022, which the re-base raised by \SI{11.07}{\mega\tonne}---Table 3's flat
conversion had been reading that year's coal low---so the shortfall became an excess and
the clean conclusion with it. This release lowers the same year again---household dung on
the survey, wastewater nitrous oxide on the 2019 Refinement with each year's own sanitation
shares, fertiliser feedstock outside the post-meter leak base---and the excess is a
shortfall once more. Both components were
always there; what changes is which one
dominates. The completeness difference is unaffected and still \SI{5.658}{\mega\tonne};
the estimation difference is again the smaller term, and the submission's own coal energy
descends from the same HDIP series the re-base moved away from, so the two accounts are not
independent on exactly the category that first reversed the sign.

Two of the apparent order-of-magnitude disagreements are one disagreement. On landfill methane
this inventory is \num{0.49} of the submission; on open burning of waste it is \num{17.74}
times it. Summed over the solid stream the ratio is \num{0.68}, so what the accounts mainly
disagree about is the \emph{fate split} of the same municipal waste---how much is collected to
a disposal site against how much burns where it falls. That is a question about collection
statistics rather than about an emission factor, and neither individual ratio should be quoted
as a level error.

The managed-soils row is the one that should change a reader's mind, because it is the second
independent account to say the same thing. The submission puts direct and indirect
managed-soil nitrous oxide at \SI{58.05}{\mega\tonne} against this inventory's
\SI{25.77}{\mega\tonne}; FAOSTAT, built from different inputs, puts the same quantity at
\SI{58.50}{} against \SI{28.53}{\mega\tonne} for the reference year. Two accounts, one
official and one international, land at about twice this inventory, and the cause this paper
already publishes---the 2019 Refinement's dry-climate \textrm{EF}$_1$ against the 2006
aggregate value---predicts exactly that. A second comparator agreeing with the first is worth
more than either alone, and it moves the dry-climate choice from a defensible assumption to
the single most consequential open question in the agricultural account.

The comparison runs both ways, which a completeness argument has to. The submission reports
the metal industry as zero, where this inventory estimates iron and steel; and it reports
substitutes for ozone-depleting substances as zero, where this inventory declares the
category NE and carries an external bound beside it. A typed absence with a bound states more
than a reported zero.

One entry in that list has to be withdrawn, and the withdrawal is the more useful result. The
submission carries no biomass-burning \emph{line}, and reading that as an absence---as an
earlier version of this crosswalk did---makes this inventory look like the more complete
account. It is not. The submission folds residue burning into its land block: cropland's
non-CO\textsubscript{2} half is \SI{3.16}{\mega\tonne} beside \SI{2.50}{\mega\tonne} of
land-use CO\textsubscript{2} in the same row, and the agriculture Common Reporting Table
carries crop residue burning as a line of its own---the same quantity filed twice by one
submission, which is also why a land layer built here cannot simply adopt both. Against the
folded figure this inventory's field burning is \num{0.18}, not a surplus. The
exclusion of the land block is therefore not as clean as scope makes it sound:
\SI{3.47}{\mega\tonne} of its \SI{34.82}{\mega\tonne} is CH\textsubscript{4} and
N\textsubscript{2}O, which for a land account is fire, and about \SI{3.5}{\mega\tonne} of
what the scope exclusion withholds is comparable in kind to emissions this inventory does
estimate. It is also the hazard to avoid in reverse: a land layer here that adopted the
submission's cropland figure beside this inventory's existing residue burning would count the
burning twice.

\subsection{Provincial SO\textsubscript{2}: where national agreement hides a fivefold gap}
\label{sec:benchmarks-provincial-so2}

Every comparison above is national, and a national ratio is the one statistic a spatial
inventory can pass without being right anywhere. The provincial test is
\texttt{outputs/pakemis\_air\_provincial\_external\_comparison.csv}, which places this
inventory's ADM1 allocation beside EDGAR~v8.1's gridded SO\textsubscript{2} summed by
area fraction into the same boundaries, on a fiscal blend of the two calendar years each
fiscal year spans.

At FY2021--22 the national ratio is \num{0.92}---as close to agreement as anything in
Table~\ref{tab:benchmarks}---and it is not informative. Underneath it Punjab is \num{1.25},
Khyber Pakhtunkhwa \num{1.67}, Balochistan \num{2.05} and Sindh \num{0.18}: a spread of an
order of magnitude, cancelling to within four percent of unity. A compilation can be wrong in
every province and right in the sum, and this one very nearly is.

The disagreement is not diffuse. Decomposing by stream puts it almost entirely in electricity
generation, where this inventory books \SI{167.0}{\giga\gram} against EDGAR's
\SI{409.3}{\giga\gram}, and within that in Sindh, where it books \SI{4.5}{\giga\gram} against
\SI{196.9}{\giga\gram}. That is not a defect in the allocation but a declared absence appearing
where it was declared: power-sector SO\textsubscript{2} \emph{as actually emitted} is NE in this
release, pending a control-device inventory, and what is booked is the oil-and-diesel admission
register alone---no coal. The coal fleet that EDGAR is seeing is in Sindh.

Read against the release's own screening the comparison resolves rather than merely disagrees.
Adding the uncontrolled coal screening to the booked register---the two are disjoint by
construction, the screening's oil rows covering only the remainder the register does not
reach---gives \SI{362.2}{\giga\gram}; assuming flue-gas desulphurisation at the screening's
controlled case gives \SI{202.0}{\giga\gram}. EDGAR's \SI{409.3}{\giga\gram} falls inside that
bracket, at \SI{129}{\percent} of its width. An independent compilation therefore places
Pakistani coal generation near the uncontrolled end of this inventory's own declared range,
which is the most that can be said before the control-device inventory exists, and is more than
the national ratio could say at all.

One caveat belongs with these provincial numbers rather than after them. The coordinate ingest
resolves every located station's province by dropping its coordinate into the \ang{0.05}
grid and reading the cell, and a cell is about \SI{5.5}{\kilo\metre} across.
\texttt{outputs/pakemis\_power\_province\_boundary\_audit.csv} tests all \num{77} located
stations against their curated records and against their sources' own province fields:
\num{7} are misplaced, every one of them within \SI{3.4}{\kilo\metre} of a district line.
The clearest is the Hub site, where two stations \SI{1.11}{\kilo\metre} apart already sit in
different provinces in this repository---the oil station recorded in Balochistan, the coal
station in Sindh because its ingest copies the raster rather than the Global Energy Monitor
field beside it, which says Balochistan. At FY2021--22 the coal station alone burned
\SI{25.5}{\percent} of the coal generated in Pakistan, some \SI{6.4}{\mega\tonne} of
CO\textsubscript{2}; across the ledger years the misplaced stations carry between
\SI{1.5}{\percent} and \SI{10.3}{\percent} of plant-ledger generation, the low end being the
FY2023--24 reference, where the Hub coal station barely ran.

A finer grid would not settle this, which is why the audit corrects nothing. For four of the
seven the simplified district polygons disagree with the curated record as well---ray casting
places the Hub oil station in Karachi, both Guddu units in Dera Ghazi Khan and Tapal Energy in
Lasbela---so the error is in the boundary geometry at these sites and not only in the
rasterisation of it. Province and district are also one decision here: the district raster is
built inside the province raster, and the nine-year products assert
$\mathrm{ADM1} = \mathrm{ADM2} + \text{outside} + \text{unallocated}$, so relabelling a
station's province while its mass stays binned in the old district would break the identity
that the spatial layer exists to hold. Correcting these seven is a declared provincial re-base
with a geometry decision inside it; the audit measures what it would move, and a guard holds
the list at seven so an eighth cannot appear in silence.

The audit now publishes \emph{which kind} of conflict each station has, because the seven are
three problems and not one. For \num{2} the frame disagrees with itself---the
\SI{0.05}{\degree} province raster and the district polygon name different provinces, so no
relabelling settles it. \num{1} sits in a hole: no district polygon contains the Hub coal
station at all, \SI{0.13}{\kilo\metre} from the nearest vertex. For the remaining \num{4}
the curated and source records agree and the frame dissents. \num{0} of the seven is a
station-record error, which is the result that fixes the remedy: every conflict is in the
geometry, so a boundary revision addresses all of them and a station relabelling addresses
none. The classification is a declared vocabulary rather than a vote, because the raster
reads the province file and the containing test reads the district file---a disagreement
between them can be a file difference or a resolution difference, and nothing held says
which, so counting them as two witnesses would invent an authority neither has.
